\subsection*{Justificación de la pregunta problema}

La pregunta problema se justifica porque responde a una tensión concreta del aula de matemáticas en educación media: se espera que los estudiantes desarrollen pensamiento trigonométrico, pero en la práctica predominan actividades centradas en procedimientos y respuestas rápidas \parencite{fiallolealAcercaInvestigacionEducacion2015,delgadoEstrategiaDidacticaPara2023}. En ese escenario, la IA puede abrir nuevas posibilidades, pero solo si se integra dentro de tareas bien diseñadas que mantengan el foco en el trabajo matemático del estudiante y no en la automatización de resultados \parencite{castillodelpezoUseArtificialIntelligence2024,parrasanchezPersonalizacionRecursosPara2023}.

Además, la pregunta delimita con claridad el alcance del estudio: educación media, trigonometría escolar y desarrollo del pensamiento trigonométrico. Esta delimitación permite analizar con detalle evidencias de aprendizaje, dificultades y decisiones didácticas en un contexto real de implementación, en lugar de hacer afirmaciones generales sobre tecnología educativa.

En consecuencia, la investigación no busca demostrar que la IA, por sí sola, mejora el aprendizaje. Busca comprender bajo qué condiciones de diseño e implementación una secuencia de tareas matemáticas integradas con IA puede aportar al desarrollo del pensamiento trigonométrico.
